% accessibility/hw1_ideal_compression.tex:  accessibility settings and alternative text
% for hw1_ideal_compression.tex (common/mulhearn/accessibility.txt).  Maintained by the agents;
% hw1_ideal_compression.tex has only % accessibility/hw1_ideal_compression.tex:  accessibility settings and alternative text
% for hw1_ideal_compression.tex (common/mulhearn/accessibility.txt).  Maintained by the agents;
% hw1_ideal_compression.tex has only % accessibility/hw1_ideal_compression.tex:  accessibility settings and alternative text
% for hw1_ideal_compression.tex (common/mulhearn/accessibility.txt).  Maintained by the agents;
% hw1_ideal_compression.tex has only % accessibility/hw1_ideal_compression.tex:  accessibility settings and alternative text
% for hw1_ideal_compression.tex (common/mulhearn/accessibility.txt).  Maintained by the agents;
% hw1_ideal_compression.tex has only \input{accessibility/hw1_ideal_compression} as its first line.
% Build with LuaLaTeX.
\DocumentMetadata{lang=en-US, pdfversion=2.0, pdfstandard=ua-2, tagging=on}

\makeatletter

% --- document information (needs hyperref, loaded here so the document does not have to)
\AddToHook{begindocument/before}{%
  \RequirePackage{hyperref}%
  \hypersetup{pdftitle={PHY 112L: Compressing a 2-D Ideal Gas}, pdfauthor={Michael Mulhearn},
    pdflang=en-US, pdfdisplaydoctitle, hidelinks,
    pdfsubject={\access@statement}}}

% --- the statement:  read first by assistive technology, not visible on the page
\def\access@statement{I have attempted to bring this document into compliance with the
  accessibility requirements of the Americans with Disabilities Act (Title II) and University of
  California policy (WCAG 2.1 Level AA).  If this minimal standard is insufficient or ineffective
  for your needs, please bring the matter to my attention, so that I can provide additional
  accommodations as appropriate.}
\AddToHook{begindocument/end}{%
  \tagstructbegin{tag=Figure, alt={\access@statement}}\tagmcbegin{}\tagmcend\tagstructend}

% --- alternative text for \includegraphics, by file name
\newcommand\AltText[2]{\@namedef{access@alt@#1}{#2}}
\AddToHook{begindocument/before}{%
  \NewCommandCopy\access@includegraphics\includegraphics
  \RenewDocumentCommand\includegraphics{s O{} m}{%
    \@ifundefined{access@alt@#3}%
      {\PackageWarning{accessibility}{no alternative text for #3}%
       \access@includegraphics[#2]{#3}}%
      {\access@includegraphics[alt={\@nameuse{access@alt@#3}},#2]{#3}}}}

% --- alternative text for TikZ pictures, by source file and order:  \AltTikZ{file.tex}{n}{text}
% (in the main file \CurrentFile is empty:  the main file's own name is used;  the stored text is
% put into the alt key before \tagstructbegin, which does not expand it)
\newcommand\AltTikZ[3]{\@namedef{access@tikz@#1@#2}{#3}}
\newcount\access@depth
\def\access@tikzbegin{%
  \ifnum\access@depth=0
    \edef\access@file{\ifx\CurrentFile\@empty\jobname.tex\else\CurrentFile\fi}%
    \@ifundefined{access@tikzn@\access@file}%
      {\global\@namedef{access@tikzn@\access@file}{0}}{}%
    \xdef\access@n{\the\numexpr\@nameuse{access@tikzn@\access@file}+1}%
    \global\expandafter\let\csname access@tikzn@\access@file\endcsname\access@n
    \@ifundefined{access@tikz@\access@file @\access@n}%
      {\PackageWarning{accessibility}%
         {no alternative text for TikZ picture \access@n\space in \access@file}%
       \global\let\access@tagged\@false}%
      {\edef\access@tmp{\noexpand\tagstructbegin{tag=Figure, alt={\unexpanded
         \expandafter\expandafter\expandafter
         {\csname access@tikz@\access@file @\access@n\endcsname}}}}%
       \access@tmp\tagmcbegin{}%
       \global\let\access@tagged\@true}%
  \fi
  \global\advance\access@depth 1 }
\def\access@tikzend{%
  \global\advance\access@depth -1
  \ifnum\access@depth=0 \ifx\access@tagged\@true \tagmcend\tagstructend\fi\fi}
\def\@true{true}\def\@false{false}
\AddToHook{env/tikzpicture/before}{\access@tikzbegin}
\AddToHook{env/tikzpicture/after}{\access@tikzend}

% --- code listings:  listings (pythonhighlight) cannot be tagged yet (TeX Live 2026:
% unbalanced paragraph hooks), so the python environment becomes a framed fancyvrb Verbatim,
% same font and frame, without colours
\AddToHook{begindocument/before}{%
  \RequirePackage{fancyvrb}%
  \RenewDocumentEnvironment{python}{}%
    {\VerbatimEnvironment
     \begin{Verbatim}[fontsize=\footnotesize, frame=single, rulecolor=\color{black!40}]}%
    {\end{Verbatim}}}

\makeatother

% --- alternative text
\AltText{plots/mbvelocity.pdf}{Graph of the probability density, in units of 1 over v zero, against
  one velocity component from minus 5 to 5 v zero, for three temperatures:  three Gaussians centred
  on zero.  At T = 1 T zero (solid line) the curve is narrowest and tallest, peaking at 0.40;  at T
  = 2 T zero (dashed) it peaks at 0.28;  at T = 5 T zero (dotted) it is widest and lowest, peaking
  at 0.18.}
\AltText{plots/vinit.pdf}{Histograms of the number of molecules against velocity, for vx (circles)
  and vy (squares), just after initialization, at T = 3.0 T zero:  both are flat at about 830
  molecules per bin between minus 3 and 3 v zero and empty outside, unlike the dashed
  Maxwell-Boltzmann curve, a Gaussian peaking near 1150 molecules at zero velocity.}
\AltText{plots/vself.pdf}{The same histograms of vx (circles) and vy (squares) after 100,000
  collisions of the gas with itself, still at T = 3.0 T zero:  both now follow the dashed
  Maxwell-Boltzmann curve, peaking near 1150 molecules at zero velocity and falling to a few dozen
  at plus and minus 4.5 v zero.}
\AltTikZ{hw1_ideal_compression.tex}{1}{Drawing of a two-dimensional cylinder of width a holding
  scattered gas molecules (dots), closed at the top by a plunger with a handle.  A vertical y axis
  beside it is marked at 0, 2a and 4a;  the plunger sits at y = 4a.}
\AltTikZ{hw1_ideal_compression.tex}{2}{Diagram of a collision in the centre-of-mass frame:  molecule
  a comes in from the left with velocity u and molecule b from the right with velocity minus u,
  meeting at the centre.  They leave along a line turned by the angle theta from the incoming line,
  a with velocity w (up and to the right) and b with velocity minus w (down and to the left).}
\AltTikZ{hw1_ideal_compression.tex}{3}{Diagram of a molecule hitting the plunger:  molecule a moves
  up and to the right with velocity u, strikes the underside of the plunger, which moves downward
  with velocity minus v, and recoils down and to the right with velocity w.}
 as its first line.
% Build with LuaLaTeX.
\DocumentMetadata{lang=en-US, pdfversion=2.0, pdfstandard=ua-2, tagging=on}

\makeatletter

% --- document information (needs hyperref, loaded here so the document does not have to)
\AddToHook{begindocument/before}{%
  \RequirePackage{hyperref}%
  \hypersetup{pdftitle={PHY 112L: Compressing a 2-D Ideal Gas}, pdfauthor={Michael Mulhearn},
    pdflang=en-US, pdfdisplaydoctitle, hidelinks,
    pdfsubject={\access@statement}}}

% --- the statement:  read first by assistive technology, not visible on the page
\def\access@statement{I have attempted to bring this document into compliance with the
  accessibility requirements of the Americans with Disabilities Act (Title II) and University of
  California policy (WCAG 2.1 Level AA).  If this minimal standard is insufficient or ineffective
  for your needs, please bring the matter to my attention, so that I can provide additional
  accommodations as appropriate.}
\AddToHook{begindocument/end}{%
  \tagstructbegin{tag=Figure, alt={\access@statement}}\tagmcbegin{}\tagmcend\tagstructend}

% --- alternative text for \includegraphics, by file name
\newcommand\AltText[2]{\@namedef{access@alt@#1}{#2}}
\AddToHook{begindocument/before}{%
  \NewCommandCopy\access@includegraphics\includegraphics
  \RenewDocumentCommand\includegraphics{s O{} m}{%
    \@ifundefined{access@alt@#3}%
      {\PackageWarning{accessibility}{no alternative text for #3}%
       \access@includegraphics[#2]{#3}}%
      {\access@includegraphics[alt={\@nameuse{access@alt@#3}},#2]{#3}}}}

% --- alternative text for TikZ pictures, by source file and order:  \AltTikZ{file.tex}{n}{text}
% (in the main file \CurrentFile is empty:  the main file's own name is used;  the stored text is
% put into the alt key before \tagstructbegin, which does not expand it)
\newcommand\AltTikZ[3]{\@namedef{access@tikz@#1@#2}{#3}}
\newcount\access@depth
\def\access@tikzbegin{%
  \ifnum\access@depth=0
    \edef\access@file{\ifx\CurrentFile\@empty\jobname.tex\else\CurrentFile\fi}%
    \@ifundefined{access@tikzn@\access@file}%
      {\global\@namedef{access@tikzn@\access@file}{0}}{}%
    \xdef\access@n{\the\numexpr\@nameuse{access@tikzn@\access@file}+1}%
    \global\expandafter\let\csname access@tikzn@\access@file\endcsname\access@n
    \@ifundefined{access@tikz@\access@file @\access@n}%
      {\PackageWarning{accessibility}%
         {no alternative text for TikZ picture \access@n\space in \access@file}%
       \global\let\access@tagged\@false}%
      {\edef\access@tmp{\noexpand\tagstructbegin{tag=Figure, alt={\unexpanded
         \expandafter\expandafter\expandafter
         {\csname access@tikz@\access@file @\access@n\endcsname}}}}%
       \access@tmp\tagmcbegin{}%
       \global\let\access@tagged\@true}%
  \fi
  \global\advance\access@depth 1 }
\def\access@tikzend{%
  \global\advance\access@depth -1
  \ifnum\access@depth=0 \ifx\access@tagged\@true \tagmcend\tagstructend\fi\fi}
\def\@true{true}\def\@false{false}
\AddToHook{env/tikzpicture/before}{\access@tikzbegin}
\AddToHook{env/tikzpicture/after}{\access@tikzend}

% --- code listings:  listings (pythonhighlight) cannot be tagged yet (TeX Live 2026:
% unbalanced paragraph hooks), so the python environment becomes a framed fancyvrb Verbatim,
% same font and frame, without colours
\AddToHook{begindocument/before}{%
  \RequirePackage{fancyvrb}%
  \RenewDocumentEnvironment{python}{}%
    {\VerbatimEnvironment
     \begin{Verbatim}[fontsize=\footnotesize, frame=single, rulecolor=\color{black!40}]}%
    {\end{Verbatim}}}

\makeatother

% --- alternative text
\AltText{plots/mbvelocity.pdf}{Graph of the probability density, in units of 1 over v zero, against
  one velocity component from minus 5 to 5 v zero, for three temperatures:  three Gaussians centred
  on zero.  At T = 1 T zero (solid line) the curve is narrowest and tallest, peaking at 0.40;  at T
  = 2 T zero (dashed) it peaks at 0.28;  at T = 5 T zero (dotted) it is widest and lowest, peaking
  at 0.18.}
\AltText{plots/vinit.pdf}{Histograms of the number of molecules against velocity, for vx (circles)
  and vy (squares), just after initialization, at T = 3.0 T zero:  both are flat at about 830
  molecules per bin between minus 3 and 3 v zero and empty outside, unlike the dashed
  Maxwell-Boltzmann curve, a Gaussian peaking near 1150 molecules at zero velocity.}
\AltText{plots/vself.pdf}{The same histograms of vx (circles) and vy (squares) after 100,000
  collisions of the gas with itself, still at T = 3.0 T zero:  both now follow the dashed
  Maxwell-Boltzmann curve, peaking near 1150 molecules at zero velocity and falling to a few dozen
  at plus and minus 4.5 v zero.}
\AltTikZ{hw1_ideal_compression.tex}{1}{Drawing of a two-dimensional cylinder of width a holding
  scattered gas molecules (dots), closed at the top by a plunger with a handle.  A vertical y axis
  beside it is marked at 0, 2a and 4a;  the plunger sits at y = 4a.}
\AltTikZ{hw1_ideal_compression.tex}{2}{Diagram of a collision in the centre-of-mass frame:  molecule
  a comes in from the left with velocity u and molecule b from the right with velocity minus u,
  meeting at the centre.  They leave along a line turned by the angle theta from the incoming line,
  a with velocity w (up and to the right) and b with velocity minus w (down and to the left).}
\AltTikZ{hw1_ideal_compression.tex}{3}{Diagram of a molecule hitting the plunger:  molecule a moves
  up and to the right with velocity u, strikes the underside of the plunger, which moves downward
  with velocity minus v, and recoils down and to the right with velocity w.}
 as its first line.
% Build with LuaLaTeX.
\DocumentMetadata{lang=en-US, pdfversion=2.0, pdfstandard=ua-2, tagging=on}

\makeatletter

% --- document information (needs hyperref, loaded here so the document does not have to)
\AddToHook{begindocument/before}{%
  \RequirePackage{hyperref}%
  \hypersetup{pdftitle={PHY 112L: Compressing a 2-D Ideal Gas}, pdfauthor={Michael Mulhearn},
    pdflang=en-US, pdfdisplaydoctitle, hidelinks,
    pdfsubject={\access@statement}}}

% --- the statement:  read first by assistive technology, not visible on the page
\def\access@statement{I have attempted to bring this document into compliance with the
  accessibility requirements of the Americans with Disabilities Act (Title II) and University of
  California policy (WCAG 2.1 Level AA).  If this minimal standard is insufficient or ineffective
  for your needs, please bring the matter to my attention, so that I can provide additional
  accommodations as appropriate.}
\AddToHook{begindocument/end}{%
  \tagstructbegin{tag=Figure, alt={\access@statement}}\tagmcbegin{}\tagmcend\tagstructend}

% --- alternative text for \includegraphics, by file name
\newcommand\AltText[2]{\@namedef{access@alt@#1}{#2}}
\AddToHook{begindocument/before}{%
  \NewCommandCopy\access@includegraphics\includegraphics
  \RenewDocumentCommand\includegraphics{s O{} m}{%
    \@ifundefined{access@alt@#3}%
      {\PackageWarning{accessibility}{no alternative text for #3}%
       \access@includegraphics[#2]{#3}}%
      {\access@includegraphics[alt={\@nameuse{access@alt@#3}},#2]{#3}}}}

% --- alternative text for TikZ pictures, by source file and order:  \AltTikZ{file.tex}{n}{text}
% (in the main file \CurrentFile is empty:  the main file's own name is used;  the stored text is
% put into the alt key before \tagstructbegin, which does not expand it)
\newcommand\AltTikZ[3]{\@namedef{access@tikz@#1@#2}{#3}}
\newcount\access@depth
\def\access@tikzbegin{%
  \ifnum\access@depth=0
    \edef\access@file{\ifx\CurrentFile\@empty\jobname.tex\else\CurrentFile\fi}%
    \@ifundefined{access@tikzn@\access@file}%
      {\global\@namedef{access@tikzn@\access@file}{0}}{}%
    \xdef\access@n{\the\numexpr\@nameuse{access@tikzn@\access@file}+1}%
    \global\expandafter\let\csname access@tikzn@\access@file\endcsname\access@n
    \@ifundefined{access@tikz@\access@file @\access@n}%
      {\PackageWarning{accessibility}%
         {no alternative text for TikZ picture \access@n\space in \access@file}%
       \global\let\access@tagged\@false}%
      {\edef\access@tmp{\noexpand\tagstructbegin{tag=Figure, alt={\unexpanded
         \expandafter\expandafter\expandafter
         {\csname access@tikz@\access@file @\access@n\endcsname}}}}%
       \access@tmp\tagmcbegin{}%
       \global\let\access@tagged\@true}%
  \fi
  \global\advance\access@depth 1 }
\def\access@tikzend{%
  \global\advance\access@depth -1
  \ifnum\access@depth=0 \ifx\access@tagged\@true \tagmcend\tagstructend\fi\fi}
\def\@true{true}\def\@false{false}
\AddToHook{env/tikzpicture/before}{\access@tikzbegin}
\AddToHook{env/tikzpicture/after}{\access@tikzend}

% --- code listings:  listings (pythonhighlight) cannot be tagged yet (TeX Live 2026:
% unbalanced paragraph hooks), so the python environment becomes a framed fancyvrb Verbatim,
% same font and frame, without colours
\AddToHook{begindocument/before}{%
  \RequirePackage{fancyvrb}%
  \RenewDocumentEnvironment{python}{}%
    {\VerbatimEnvironment
     \begin{Verbatim}[fontsize=\footnotesize, frame=single, rulecolor=\color{black!40}]}%
    {\end{Verbatim}}}

\makeatother

% --- alternative text
\AltText{plots/mbvelocity.pdf}{Graph of the probability density, in units of 1 over v zero, against
  one velocity component from minus 5 to 5 v zero, for three temperatures:  three Gaussians centred
  on zero.  At T = 1 T zero (solid line) the curve is narrowest and tallest, peaking at 0.40;  at T
  = 2 T zero (dashed) it peaks at 0.28;  at T = 5 T zero (dotted) it is widest and lowest, peaking
  at 0.18.}
\AltText{plots/vinit.pdf}{Histograms of the number of molecules against velocity, for vx (circles)
  and vy (squares), just after initialization, at T = 3.0 T zero:  both are flat at about 830
  molecules per bin between minus 3 and 3 v zero and empty outside, unlike the dashed
  Maxwell-Boltzmann curve, a Gaussian peaking near 1150 molecules at zero velocity.}
\AltText{plots/vself.pdf}{The same histograms of vx (circles) and vy (squares) after 100,000
  collisions of the gas with itself, still at T = 3.0 T zero:  both now follow the dashed
  Maxwell-Boltzmann curve, peaking near 1150 molecules at zero velocity and falling to a few dozen
  at plus and minus 4.5 v zero.}
\AltTikZ{hw1_ideal_compression.tex}{1}{Drawing of a two-dimensional cylinder of width a holding
  scattered gas molecules (dots), closed at the top by a plunger with a handle.  A vertical y axis
  beside it is marked at 0, 2a and 4a;  the plunger sits at y = 4a.}
\AltTikZ{hw1_ideal_compression.tex}{2}{Diagram of a collision in the centre-of-mass frame:  molecule
  a comes in from the left with velocity u and molecule b from the right with velocity minus u,
  meeting at the centre.  They leave along a line turned by the angle theta from the incoming line,
  a with velocity w (up and to the right) and b with velocity minus w (down and to the left).}
\AltTikZ{hw1_ideal_compression.tex}{3}{Diagram of a molecule hitting the plunger:  molecule a moves
  up and to the right with velocity u, strikes the underside of the plunger, which moves downward
  with velocity minus v, and recoils down and to the right with velocity w.}
 as its first line.
% Build with LuaLaTeX.
\DocumentMetadata{lang=en-US, pdfversion=2.0, pdfstandard=ua-2, tagging=on}

\makeatletter

% --- document information (needs hyperref, loaded here so the document does not have to)
\AddToHook{begindocument/before}{%
  \RequirePackage{hyperref}%
  \hypersetup{pdftitle={PHY 112L: Compressing a 2-D Ideal Gas}, pdfauthor={Michael Mulhearn},
    pdflang=en-US, pdfdisplaydoctitle, hidelinks,
    pdfsubject={\access@statement}}}

% --- the statement:  read first by assistive technology, not visible on the page
\def\access@statement{I have attempted to bring this document into compliance with the
  accessibility requirements of the Americans with Disabilities Act (Title II) and University of
  California policy (WCAG 2.1 Level AA).  If this minimal standard is insufficient or ineffective
  for your needs, please bring the matter to my attention, so that I can provide additional
  accommodations as appropriate.}
\AddToHook{begindocument/end}{%
  \tagstructbegin{tag=Figure, alt={\access@statement}}\tagmcbegin{}\tagmcend\tagstructend}

% --- alternative text for \includegraphics, by file name
\newcommand\AltText[2]{\@namedef{access@alt@#1}{#2}}
\AddToHook{begindocument/before}{%
  \NewCommandCopy\access@includegraphics\includegraphics
  \RenewDocumentCommand\includegraphics{s O{} m}{%
    \@ifundefined{access@alt@#3}%
      {\PackageWarning{accessibility}{no alternative text for #3}%
       \access@includegraphics[#2]{#3}}%
      {\access@includegraphics[alt={\@nameuse{access@alt@#3}},#2]{#3}}}}

% --- alternative text for TikZ pictures, by source file and order:  \AltTikZ{file.tex}{n}{text}
% (in the main file \CurrentFile is empty:  the main file's own name is used;  the stored text is
% put into the alt key before \tagstructbegin, which does not expand it)
\newcommand\AltTikZ[3]{\@namedef{access@tikz@#1@#2}{#3}}
\newcount\access@depth
\def\access@tikzbegin{%
  \ifnum\access@depth=0
    \edef\access@file{\ifx\CurrentFile\@empty\jobname.tex\else\CurrentFile\fi}%
    \@ifundefined{access@tikzn@\access@file}%
      {\global\@namedef{access@tikzn@\access@file}{0}}{}%
    \xdef\access@n{\the\numexpr\@nameuse{access@tikzn@\access@file}+1}%
    \global\expandafter\let\csname access@tikzn@\access@file\endcsname\access@n
    \@ifundefined{access@tikz@\access@file @\access@n}%
      {\PackageWarning{accessibility}%
         {no alternative text for TikZ picture \access@n\space in \access@file}%
       \global\let\access@tagged\@false}%
      {\edef\access@tmp{\noexpand\tagstructbegin{tag=Figure, alt={\unexpanded
         \expandafter\expandafter\expandafter
         {\csname access@tikz@\access@file @\access@n\endcsname}}}}%
       \access@tmp\tagmcbegin{}%
       \global\let\access@tagged\@true}%
  \fi
  \global\advance\access@depth 1 }
\def\access@tikzend{%
  \global\advance\access@depth -1
  \ifnum\access@depth=0 \ifx\access@tagged\@true \tagmcend\tagstructend\fi\fi}
\def\@true{true}\def\@false{false}
\AddToHook{env/tikzpicture/before}{\access@tikzbegin}
\AddToHook{env/tikzpicture/after}{\access@tikzend}

% --- code listings:  listings (pythonhighlight) cannot be tagged yet (TeX Live 2026:
% unbalanced paragraph hooks), so the python environment becomes a framed fancyvrb Verbatim,
% same font and frame, without colours
\AddToHook{begindocument/before}{%
  \RequirePackage{fancyvrb}%
  \RenewDocumentEnvironment{python}{}%
    {\VerbatimEnvironment
     \begin{Verbatim}[fontsize=\footnotesize, frame=single, rulecolor=\color{black!40}]}%
    {\end{Verbatim}}}

\makeatother

% --- alternative text
\AltText{plots/mbvelocity.pdf}{Graph of the probability density, in units of 1 over v zero, against
  one velocity component from minus 5 to 5 v zero, for three temperatures:  three Gaussians centred
  on zero.  At T = 1 T zero (solid line) the curve is narrowest and tallest, peaking at 0.40;  at T
  = 2 T zero (dashed) it peaks at 0.28;  at T = 5 T zero (dotted) it is widest and lowest, peaking
  at 0.18.}
\AltText{plots/vinit.pdf}{Histograms of the number of molecules against velocity, for vx (circles)
  and vy (squares), just after initialization, at T = 3.0 T zero:  both are flat at about 830
  molecules per bin between minus 3 and 3 v zero and empty outside, unlike the dashed
  Maxwell-Boltzmann curve, a Gaussian peaking near 1150 molecules at zero velocity.}
\AltText{plots/vself.pdf}{The same histograms of vx (circles) and vy (squares) after 100,000
  collisions of the gas with itself, still at T = 3.0 T zero:  both now follow the dashed
  Maxwell-Boltzmann curve, peaking near 1150 molecules at zero velocity and falling to a few dozen
  at plus and minus 4.5 v zero.}
\AltTikZ{hw1_ideal_compression.tex}{1}{Drawing of a two-dimensional cylinder of width a holding
  scattered gas molecules (dots), closed at the top by a plunger with a handle.  A vertical y axis
  beside it is marked at 0, 2a and 4a;  the plunger sits at y = 4a.}
\AltTikZ{hw1_ideal_compression.tex}{2}{Diagram of a collision in the centre-of-mass frame:  molecule
  a comes in from the left with velocity u and molecule b from the right with velocity minus u,
  meeting at the centre.  They leave along a line turned by the angle theta from the incoming line,
  a with velocity w (up and to the right) and b with velocity minus w (down and to the left).}
\AltTikZ{hw1_ideal_compression.tex}{3}{Diagram of a molecule hitting the plunger:  molecule a moves
  up and to the right with velocity u, strikes the underside of the plunger, which moves downward
  with velocity minus v, and recoils down and to the right with velocity w.}
